\section{Proxy discrimination and interventions}
\label{sec:proxies}

We now turn to an important aspect of our framework. Determining causal effects
in general requires modeling interventions. Interventions on deeply rooted
individual properties such as \emph{gender} or \emph{race} are notoriously
difficult to conceptualize---especially at an individual level, and impossible
to perform in a randomized trial. VanderWeele et al.~\cite{VanderWeele2014}
discuss the problem comprehensively in an epidemiological setting. From a
machine learning perspective, it thus makes sense to separate the protected
attribute~$A$ from its potential \emph{proxies}, such as name, visual features,
languages spoken at home, etc. Intervention based on proxy variables poses a
more manageable problem. By deciding on a suitable proxy we can find an adequate
mounting point for determining and removing its influence on the prediction.
Moreover, in practice we are often limited to imperfect measurements of~$A$ in
any case, making the distinction between root concept and proxy prudent.

As was the case with resolving variables, a \emph{proxy} is a priori nothing
more than a descendant of~$A$ in the causal graph that we choose to label as a
proxy. Nevertheless in reality we envision the proxy to be a clearly defined
observable quantity that is significantly correlated with~$A,$ yet in our view
should not affect the prediction.

\begin{definition}[Potential proxy discrimination]\label{def:proxy}
  A variable~$V$ in a causal graph exhibits \emph{potential proxy
  discrimination}, if there exists a directed path from~$A$ to~$V$ that is
  blocked by a proxy variable and~$V$ itself is not a proxy.
\end{definition}

Potential proxy discrimination articulates a causal criterion that is in a sense
dual to unresolved discrimination. From the \emph{benevolent viewpoint}, we \emph{allow} any path from~$A$
to~$\hR$ unless it passes through a proxy variable, which we consider worrisome.
This viewpoint acknowledges the fact that the influence of~$A$ on the graph may
be complex and it can be too restraining to rule out all but a few designated
features. In practice, as with unresolved discrimination, we can naively build
an unconstrained predictor based only on those features that do not exhibit
potential proxy discrimination. Then we must not provide~$P$ as input to~$R;$
unawareness, \ie{} excluding~$P$ from the inputs of~$R$, suffices. However, by
granting~$R$ access to~$P$, we can carefully tune the function~$R(P, X)$ to
cancel the implicit influence of~$P$ on features~$X$ that exhibit potential
proxy discrimination by the explicit dependence on $P$. Due to this possible
cancellation of paths, we called the path based criterion \emph{potential} proxy
discrimination. When building predictors that exhibit no \emph{overall proxy
discrimination}, we precisely aim for such a cancellation.

Fortunately, this idea can be conveniently expressed by an \emph{intervention}
on~$P$, which is denoted by~$do(P=p)$~\cite{Pearl2009}. Visually, intervening
on~$P$ amounts to removing all incoming arrows of~$P$ in the graph;
algebraically, it consists of replacing the structural equation of $P$ by $P=p$,
\ie{} we put point mass on the value~$p$.

\begin{definition}[Proxy discrimination]\label{def:overalldisc}
  A predictor~$\hR$ exhibits no \emph{proxy discrimination} based on a proxy~$P$
  if for all~$p, p'$
  \begin{equation}\label{eq:proxy-eq}
    \bP(\hR \given do(P=p) ) = \bP(\hR \given do(P=p') ) \eqp
  \end{equation}
\end{definition}

The interventional characterization of proxy discrimination leads to a simple
procedure to remove it in causal graphs that we will turn to in the next
section. It also leads to several natural variants of the definition that we
discuss in Section~\ref{sec:variants}. We remark that
Equation~\eqref{eq:proxy-eq} is an equality of probabilities in the
``do-calculus'' that cannot in general be inferred by an observational method,
because it depends on an underlying causal graph, see the discussion
in~\cite{Pearl2009}. However, in some cases, we do not need to resort to
interventions to avoid proxy discrimination.

\begin{proposition}\label{pro:unawareness}
  If there is no directed path from a proxy to a feature, unawareness avoids
  proxy discrimination.
\end{proposition}
